\emph{Preliminaries.} By Theorem \claimref{synthesis.signed-directions-and-the-integer-clifford-algebra} the integer algebra is free on the blades; the adjoint, observables and (P) are those of Theorem \claimref{synthesis.chsh-bound-for-two-subsystems-of-a-clifford-algebra-by-a-rational-certificate}, whose preliminaries give $\langle x^{\dagger}\rangle_{0} = \langle x\rangle_{0}$, $\langle ab\rangle_{0} = \langle ba\rangle_{0}$, that a self-adjoint blade has square $1$, and that the scalar part of a product of two blades is zero unless they are equal up to sign. \emph{The configuration.} $p_{1} = f_{1}e_{1}$ is self-adjoint of square $1$, anticommutes with $e_{1}$, $f_{1}$ and commutes with $e_{2}$, $f_{2}$; hence $e_{2}p_{1}$ has square $e_{2}^{2}p_{1}^{2} = 1$, is self-adjoint, commutes with $e_{1}$ ($e_{1}e_{2}p_{1} = -e_{2}e_{1}p_{1} = e_{2}p_{1}e_{1}$) and with $p_{1}$, and anticommutes with $p_{2}$ as $e_{2}$ does; this is the twisted generator of Theorem \claimref{tower.commuting-copies}. So the four blades are a configuration with anticommuting pairs, and clauses (iv), (v) of Theorem \claimref{synthesis.chsh-bound-for-two-subsystems-of-a-clifford-algebra-by-a-rational-certificate} give $C^{2} = 4 + 4K$, $K^{\dagger} = K$, $K^{2} = 1$, $\langle K\rangle_{0} = 0$, and for $\lambda$: $\lambda(1) = 1$, $\lambda(x^{\dagger}) = \lambda(x)$, $\lambda(C) = 4$, $\lambda((C - t)(C - t)^{\dagger}) = t^{2} - 8t + 4$.

\emph{(i).} The four products $A_{i}B_{j}$ are, up to sign, the blades $p_{1}p_{2}$, $p_{1}e_{2}p_{1} = e_{2}$, $e_{1}p_{2}$, $e_{1}e_{2}p_{1}$: distinct, not scalars, self-adjoint of square $1$, and none equals $\pm$ an $A_{i}$, a $B_{j}$ or $K$ (the check lists them). For a blade $b$ of square $1$, $\lambda(b)$ is the coefficient of $b$ in $1 + C$: this gives the values $1, 1, 1, -1$ on the products and $0$ on the $A_{i}$, the $B_{j}$ and $K$. Expanding, $\lambda(\tfrac14(1 + aA_{i})(1 + bB_{j})) = \tfrac14(1 + ab\,\lambda(A_{i}B_{j}))$, which is $\tfrac12$ or $0$; summing over $b$ gives $\tfrac12$, independent of $j$, and likewise over $a$. \emph{One setting pair.} $A_{i}$ and $B_{j}$ are commuting self-adjoint elements of square $1$, so the algebra they generate is spanned by $1, A_{i}, B_{j}, A_{i}B_{j}$, equivalently by the four idempotents $E_{ab} = \tfrac14(1 + aA_{i})(1 + bB_{j})$, which are self-adjoint, mutually orthogonal and sum to $1$. For $x = \sum x_{ab}E_{ab}$, $xx^{\dagger} = \sum x_{ab}^{2}E_{ab}$ and $\lambda(xx^{\dagger}) = \sum x_{ab}^{2}\lambda(E_{ab}) \ge 0$.

\emph{(ii).} The Gram matrix $\lambda(e_{S}e_{T}^{\dagger})$ on the sixteen blades is a rational symmetric matrix (symmetric because $\lambda(x^{\dagger}) = \lambda(x)$); exact congruence diagonalisation gives $12$ positive, $4$ negative and no zero diagonal entries (the check), and by the law of inertia this is the largest dimension of a subspace on which the form is negative definite. The quadratic $t^{2} - 8t + 4$ has roots $4 \pm 2\sqrt{3}$. For $h = A_{1}(B_{1} + B_{2})$: $h^{\dagger} = h$, $h^{2} = (B_{1} + B_{2})^{2} = 2 + B_{1}B_{2} + B_{2}B_{1} = 2$, $\lambda(h) = 2$, so $\lambda((h - 2)^{2}) = 2 - 8 + 4 = -2$. An element $x$ with $\lambda(xx^{\dagger}) < 0$ cannot lie in the algebra of one setting pair, by (i).

\emph{(iii).} $\varphi((C - t)^{2}) = \varphi(C^{2}) - 2t\varphi(C) + t^{2}$, which at $t = \varphi(C)$ is $\varphi(C^{2}) - \varphi(C)^{2}$; so the first condition is $\varphi(C)^{2} \le \varphi(C^{2})$. $(1 - K)^{2} = 2 - 2K$, so the second gives $\varphi(K) \le 1$ and $\varphi(C^{2}) = 4 + 4\varphi(K) \le 8$. $C$ contains all four products $A_{i}B_{j}$, and $K$ is, up to sign, $A_{1}B_{1} \cdot A_{2}B_{2}$. For $\lambda$: $\lambda((1 - K)^{2}) = 2 - 2\lambda(K) = 2$, and $\lambda((C - 4)^{2}) = -12$. Since $\lambda$ is a normalised linear functional on the algebra with $\lambda(x^{\dagger}) = \lambda(x)$, satisfies every identity of elements, and has $\lambda(C)^{2} = 16 > 8$, no argument from the algebra, linearity, normalisation and reality alone gives the bound; and since a sum of elements $xx^{\dagger}$ would make $\langle(1 + C)\,\cdot\,\rangle_{0}$ non-negative on every $yy^{\dagger}$ (as $\langle xx^{\dagger}yy^{\dagger}\rangle_{0} = \langle (y^{\dagger}x)^{\dagger}(y^{\dagger}x)\rangle_{0} \ge 0$), $1 + C$ is not such a sum.
